\section{Исследование смещения центральной частоты настройки}

В данном разделе исследуется влияние расстройки центральной частоты несущего колебания относительно номинального значения на качество демодуляции и спектральные характеристики тракта приема.

\subsection{Результат эксперимента}
Для анализа влияния знака и величины частотной расстройки проводилось смещение центральной частоты настройки на величину $\Delta f = \pm 4\text{ МГц}$ относительно номинального значения несущей $f_0$. Полученные спектральные распределения и отклики квадратурных каналов представлены на рис.~\ref{fig:freq_shift_minus4} и рис.~\ref{fig:freq_shift_plus4}.

\begin{figure}[htbp]
    \centering
    \begin{minipage}[b]{0.48\textwidth}
        \centering
        \includegraphics[width=\textwidth]{fig2.png}
        \caption{Отклик тракта при отрицательной расстройке ($\Delta f = -4\text{ МГц}$)}
        \label{fig:freq_shift_minus4}
    \end{minipage}
    \hfill
    \begin{minipage}[b]{0.48\textwidth}
        \centering
        \includegraphics[width=\textwidth]{fig3.png}
        \caption{Отклик тракта при положительной расстройке ($\Delta f = +0.41\text{ МГц}$)}
        \label{fig:freq_shift_plus4}
    \end{minipage}
\end{figure}

При установке $\text{Center Freq} = 995.98\text{ МГц}$ спектральный пик входного сигнала смещается вправо относительно центра экрана. При $\text{Center Freq} = 1000.39\text{ МГц}$ пик смещается влево.

Поскольку частота сигнала генератора фиксирована ($1.0\text{ ГГц}$), при увеличении центральной частоты анализатора Center Freq центр шкалы смещается в область более высоких частот, поэтому сам пик сигнала на экране визуально смещается влево. И наоборот: при уменьшении Center Freq пик смещается вправо.
\newpage
